\subsection{Relative central classes of product pairs}
\label{sec:cm-cross-prime-central}

Let \(\{(C_i,D_i)\}_{i\in I}\) be abelian profinite pairs, with
\(D_i\) closed in \(C_i\), and put
\[
 C=\prod_{i\in I}C_i,\qquad D=\prod_{i\in I}D_i,
 \qquad k=\mathbf F_2.
\]
The coefficient action of all these groups on \(k\) is trivial.
Write
\[
 W_i=\operatorname{Hom}_{\rm cont}(C_i,k),\quad
 V_i=\operatorname{Hom}_{\rm cont}(D_i,k),\quad
 U_i=\ker(W_i\longrightarrow V_i).
\]
For distinct indices let
\[
 K_{ij}=\ker\bigl(W_i\otimes_k W_j
                  \longrightarrow V_i\otimes_k V_j\bigr).
\]
Equivalently, \(K_{ij}\) is the space of continuous bilinear maps
\(C_i\times C_j\to k\) that vanish on \(D_i\times D_j\).
Every such bilinear map factors through finite elementary abelian
quotients, which justifies its identification with a tensor of the
continuous character spaces. Linear algebra over the field \(k\)
gives
\begin{equation}
 K_{ij}=U_i\otimes W_j+W_i\otimes U_j
               \quad\hbox{inside }W_i\otimes W_j.       \label{eq:cm-mixed-central-module}
\end{equation}
No surjectivity of \(W_i\to V_i\) is asserted or needed.

Fix an ordering of \(I\) to write each unordered pair once.

\begin{proposition}[Relative product decomposition]
\label{prop:cm-relative-product-cohomology}
There is a natural isomorphism
\begin{equation}
 H^2_{\rm cont}(C,D;k)
 \simeq
 \bigoplus_{i\in I}H^2_{\rm cont}(C_i,D_i;k)
 \ \oplus\!\!\bigoplus_{i<j}K_{ij}.                    \label{eq:cm-relative-product-cohomology}
\end{equation}
It is equivariant for every group acting continuously and
coordinatewise on the pairs. The first coordinates are restrictions
to individual pairs. The mixed coordinate of a relative central
extension is the commutator pairing between its two indicated
coordinate subgroups. All sums are algebraic direct sums, so a class
has finite coordinate support.
\end{proposition}

\begin{proof}
First suppose \(I\) is finite. We use the standard elementary
description of continuous relative degree--two classes as central
extensions by \(k\), equipped with a prescribed homomorphic splitting
over \(D\), up to equivalence preserving that splitting. It follows
directly from normalized continuous cocycles: a cocycle which is zero
on \(D\times D\) defines the group law on \(k\times C\), and
relative coboundaries are changes of section that are zero on \(D\).
Conversely a finite-kernel profinite extension has a continuous
set-theoretic section, and this section can be chosen to agree with
the prescribed one on the closed subgroup \(D\).

For completeness, choose an open subgroup of the extension meeting
the finite kernel trivially. Its map to its open image in \(C\)
is a homeomorphism; choosing finitely many coset representatives
constructs a continuous section. The discrepancy from the prescribed
section on \(D\) is a continuous finite-valued function. Such a
function on a closed subset of a profinite space extends to the
whole space by clopen separation, so the section can be corrected.

Let \(\widetilde C\to C\) represent a relative central class.
Because \(C\) is abelian, all commutators in \(\widetilde C\)
lie in its central kernel. Consequently the commutator of lifts of
\(x_i\in C_i\) and \(y_j\in C_j\) is independent of their
lifts and is a continuous bilinear function \(b_{ij}(x_i,y_j)\).
The prescribed splitting is a homomorphism on \(D\), so
\(b_{ij}\) vanishes on \(D_i\times D_j\). This gives the map
in \eqref{eq:cm-relative-product-cohomology}, together with the individual restrictions. The construction
respects Baer sum.

The map is surjective. An individual relative extension can be
pulled back by \(C\to C_i\), using the corresponding prescribed
splitting on \(D\); this has the stated individual component and
zero other components. For \(b\in K_{ij}\), the continuous
normalized cocycle
\[
 \alpha_b(x,y)=b(x_i,y_j)
\]
defines a relative extension because it is zero on
\(D\times D\). Its individual restrictions vanish, and its only
mixed commutator pairing is \(b\). Finite Baer sums realize all
the prescribed coordinates.

For injectivity, suppose all individual and mixed components
vanish. The restriction over each \(C_i\) then has a continuous
homomorphic section \(\sigma_i\) agreeing with the prescribed
section on \(D_i\). Vanishing of every mixed commutator says that
these sections have mutually commuting images. Their product
\[
 \sigma((x_i)_i)=\prod_i\sigma_i(x_i)
\]
is therefore a continuous homomorphic section over \(C\).
It agrees with the prescribed section on \(D\), since that
section is itself a homomorphism and agrees on each factor.
The original relative class is zero. This proves \eqref{eq:cm-relative-product-cohomology} for finite
products. Each step is coordinatewise natural, giving equivariance.

Now allow \(I\) to be infinite. In every degree, restriction of
continuous \(k\)-valued cochains from \(C\) to \(D\) is
surjective by the same clopen extension argument. Thus the kernel
of restriction is a relative cochain model. A continuous function
\(C^n\to k\) depends on finitely many coordinates: choose
cylinder neighborhoods on which it is constant, take a finite
subcover by compactness, and collect the finitely many coordinate
sets used by that subcover.

If such a cochain is zero on \(D^n\) and depends on a finite set
\(F\subset I\), the induced cochain on \(C_F^n\) is zero on
\(D_F^n\), where \(C_F=\prod_{i\in F}C_i\) and similarly
for \(D_F\). Hence the relative cochain complex is the filtered
union of the finite-coordinate relative complexes. Filtered unions
of vector spaces are exact, so its cohomology is the corresponding
filtered colimit. Pullback from a finite product is injective because
coordinate inclusion is a retraction as a map of pairs. Applying the
finite-product calculation proves \eqref{eq:cm-relative-product-cohomology}, with finite support as
stated.
\end{proof}

\subsection{Eliminating the mixed invariant summands}

Let a group \(\mathcal G\) act coordinatewise on the pairs.
Suppose that for every \(i\) there is an operator \(S_i\) induced
by an element of \(\mathcal G\), acting trivially at every
coordinate other than \(i\), such that its action on \(C_i\)
has finite order dividing \(2^{a_i}d_i\), with \(d_i\) odd.
Assume further that \(S_i-1\) acts bijectively on \(U_i\).
Here and below the action on character spaces is the dual action.

\begin{theorem}[Absence of mixed invariant central classes]
\label{thm:cm-cross-prime-central}
Under these hypotheses,
\begin{equation}
 H^2_{\rm cont}(C,D;k)^{\mathcal G}
 \simeq
 \bigoplus_{i\in I}
       H^2_{\rm cont}(C_i,D_i;k)^{\mathcal G}.         \label{eq:cm-prime-central-invariants}
\end{equation}
The isomorphism is restriction; its inverse is the finite sum of
coordinate pullbacks. No assertion that the individual invariant
groups vanish is part of the theorem.
\end{theorem}

\begin{proof}
Set \(T_i=S_i^{2^{a_i}}\), so its action on the \(i\)-th
coordinate has odd order dividing \(d_i\). In characteristic two,
\[
 T_i-1=(S_i-1)^{2^{a_i}}
                    \quad\hbox{on }U_i.
\]
This is invertible by hypothesis. The averaging operator
\[
 N_i=\sum_{a=0}^{d_i-1}T_i^a
\]
is the projection onto the \(T_i\)-fixed subspace: the usual
normalizing factor \(d_i^{-1}\) is one in \(k\). In particular
\(N_iU_i=0\). Its action on every other coordinate is the identity.

Let \(b\in K_{ij}^{\mathcal G}\). Since \(T_i,T_j\) come
from elements of \(\mathcal G\), invariance gives
\(N_iN_jb=b\). On the other hand \eqref{eq:cm-mixed-central-module} expresses \(b\) as
a sum from \(U_i\otimes W_j\) and \(W_i\otimes U_j\).
The commuting operator product \(N_iN_j\) kills the first
summand by \(N_iU_i=0\), and the second by \(N_jU_j=0\).
Thus \(b=0\).

The decomposition \eqref{eq:cm-relative-product-cohomology} is equivariant and every class has finite
support. Taking invariants therefore gives precisely \eqref{eq:cm-prime-central-invariants}.
\end{proof}

\subsection{The individual-prime central obstructions}

Use the notation of Proposition~\ref{prop:all-cm-exact-products}:
\[
 C=\prod_{p\in\mathcal I}C_p,\qquad
 D_0=\prod_{p\in\mathcal I}D_p,\qquad
 U_p=\operatorname{Hom}_{\rm cont}(C_p/D_p,\mathbf F_2).
\]
These are the full and arithmetic images; neither is replaced
by an ambient torus or a lower bound for the arithmetic image.

\begin{corollary}\label{cor:cm-individual-central-obstructions}
Restriction induces an isomorphism
\[
 H^2_{\rm cont}(C,D_0;\mathbf F_2)^G
 \simeq
 \bigoplus_{p\in\mathcal I}
       H^2_{\rm cont}(C_p,D_p;\mathbf F_2)^G.
\]
Every intrinsic relative central obstruction in
Theorem~\ref{thm:all-cm-central-extension} has only finitely many
nonzero individual-prime components and has no mixed component.
The corresponding character extends to a continuous character of
\(G\) vanishing on \(H\) if and only if all its individual-prime
components vanish. Such an extension is unique.
\end{corollary}

\begin{proof}
A primitive arithmetic \(p\)-inertia element induces \(s_p\)
at \(p\) and the identity at every other coordinate. Its coordinate
action has order dividing \(p-1\). By
Lemma~\ref{lem:all-cm-coordinate-module}, \(s_p-1\) is invertible
on \(U_p\). Theorem~\ref{thm:cm-cross-prime-central} therefore
applies to the exact products of
Proposition~\ref{prop:all-cm-exact-products}.

Each intrinsic relative central obstruction is invariant. Conjugation
by the arithmetic subgroup preserves its prescribed arithmetic
splitting, and the arithmetic subgroup supplies the same action on
\(C\) as \(G=R_0H\), since \(R_0\) acts trivially on the
abelian scalar image. The decomposition and finite-support assertion
follow. The extension assertion and uniqueness now follow from
Theorem~\ref{thm:all-cm-central-extension}.
\end{proof}

This result leaves the individual relative central groups to be
determined. In particular, it does not give a finite list of all
nonzero individual groups as the primes vary. At degree \(163\),
when \((C_{163},D_{163})=(P,A)\), the quaternion line of
Proposition~\ref{prop:163-relative-cohomology} is one of the
individual invariant summands. Its vanishing for full GT is not
decided by the product decomposition.
